\begin{abstract}
A terminal with low Earth orbit (LEO) satellite and terrestrial 5G paths can schedule traffic ahead of time only if it knows how each path will behave in the next seconds and how certain that knowledge is. Using 5\,Hz measurements of Starlink, OneWeb, and 5G links, we forecast the round-trip time (RTT), jitter, and loss 0.2\,s to 3\,s ahead. First, we show that the common practice of aligning ping streams by time interpolation leaks the next sample into the model input, prove that it halves the one-step persistence error on every interpolated row, and measure a \AuditPersHonePct\% understatement of the first-step error on Starlink. Second, we propose a recurrent encoder-decoder anchored on the last observation whose heads output a mean and five non-crossing quantiles per target and step, trained with a squared error plus pinball loss, and calibrated on held-out sessions with conformalized quantile regression. Against the same network with a scalar squared-error output, with three seeds each, the 95\% interval score of the RTT forecasts improves on all six links by \ChgMISWorst\% to \ChgMISBest\% (mean \ChgMISMean\%) at a measured coverage of \PICPMin{} to \PICPMax{}, and the RTT mean absolute error drops by \ChgMAEWorst\% to \ChgMAEBest\%. The RTT root-mean-square error changes by only \ChgRMSEBest\% to $+$\ChgRMSEWorst\%, and a tuned linear model is within one point of the neural baseline's skill on four links. We conclude that the value of learning lies in input-dependent uncertainty that a multipath scheduler can use as a per-path risk bound.
\end{abstract}
\begin{IEEEkeywords}
Conformal prediction, LEO satellite networks, link-state forecasting, multipath transport, OneWeb, quantile regression, Starlink, 5G.
\end{IEEEkeywords}
